% General-prime strengthening of two_adic_lower_bound.tex.
\documentclass[11pt]{article}
\usepackage{amsmath,amssymb,amsthm,hyperref}
\newtheorem{theorem}{Theorem}
\newtheorem{lemma}[theorem]{Lemma}
\newcommand{\Q}{\mathbb Q}
\newcommand{\Z}{\mathbb Z}
\begin{document}
\title{Rational interval designs: a general-prime lower bound}
\author{Research note; independently reviewed}
\date{30 September 2026}
\maketitle

\begin{theorem}
Suppose rational numbers $x_1,\ldots,x_K$ satisfy
\[
 \frac1K\sum_{i=1}^Kx_i^j=\frac1{j+1}\qquad(0\leq j\leq r).
\]
Then, for every prime $p$,
\[
 K\geq p^{\lfloor r/p\rfloor}.
\]
In particular $K\geq\max(2^{\lfloor r/2\rfloor},3^{\lfloor r/3\rfloor})$.
If all nodes are $p$-adically integral, then
$p^{\lfloor r/(p-1)\rfloor}\mid K$.
\end{theorem}

Fix a primitive $p$th root of unity $\zeta$ in an algebraic extension of
$\Q_p$, and put $\pi=\zeta-1$. Write $w$ for the valuation normalized by
$w(\pi)=1$, so $w(q)=(p-1)v_p(q)$ for rational $q$.
The extension $\Q_p(\zeta)/\Q_p$ is totally ramified of degree $p-1$.

Let $G(z)=z/\log(1+z)=\sum_{j\geq0}g_jz^j$.
Then
\[
 \frac1{G(\pi t)}
  =\sum_{j\geq0}\frac{(-1)^j\pi^j}{j+1}t^j.
\]
All these coefficients are integral for $w$. Indeed,
$j\geq(p-1)v_p(j+1)$ for $j\geq0$.
Equality holds only for $j=0,p-1$. The cyclotomic equation gives
$\pi^{p-1}/p\equiv-1\pmod\pi$, so
\[
 \frac1{G(\pi t)}\equiv1-t^{p-1}\pmod\pi.
\]
Inverting shows
\begin{equation}\label{eq:gj}
 w(g_j)\geq-j\quad(j\geq0),\qquad
 w(g_j)=-j\quad\hbox{if }p-1\mid j.
\end{equation}
Coefficientwise integration identifies
$g_j=\int_0^1\binom{x}{j}\,dx$.
If every node is in $\Z_p$, apply exactness to
$\binom{x}{s}$, where $s=(p-1)\lfloor r/(p-1)\rfloor$.
Since $\binom{x_i}{s}\in\Z_p$, \eqref{eq:gj} gives the asserted divisibility
of $K$ and hence its lower bound.

\paragraph{Historical attribution.}
The unit-indicator filter used next is classical. Mit'kin's Lemma~7
\cite{mitkin} gives the cyclotomic-trace polynomial for degrees divisible
by $p-1$, with the same congruence precision. His following remark credits
the $p=2$ binomial polynomial to Hua (1940) and records Arkhipov's related
use in the Hilbert--Kamke problem. The derivation below is included for
completeness; the filter itself is not claimed as new. The interval-integral
estimate and its application to rational equal-weight quadrature are the
specific steps studied in this note; their literature novelty remains
unverified.

For the remaining case, define the rational degree-$r$ polynomial
\[
 H_{r,p}(y)=-\frac1p\sum_{h=1}^{p-1}\sum_{j=1}^r
                   (\zeta^h-1)^j\binom yj.
\]
Its coefficients in the binomial basis are integers, since for $j\geq1$,
\[
 -\frac1p\sum_{h=1}^{p-1}(\zeta^h-1)^j
  =-\sum_{\substack{0\leq k\leq j\\p\mid k}}
                       (-1)^{j-k}\binom jk.
\]
The root-of-unity filter and the binomial theorem give, first for
nonnegative integers $y$ and then by continuity for all $y\in\Z_p$,
\begin{equation}\label{eq:filter}
 H_{r,p}(y)\equiv\boldsymbol1_{y\notin p\Z_p}
       \pmod {p^{\lfloor r/(p-1)\rfloor}}.
\end{equation}
To see the precision, every omitted term of index $j>r$ has
$w$-valuation at least $j-(p-1)$; the total is rational, so its
$p$-adic valuation is at least
$\lceil(r+1)/(p-1)-1\rceil=\lfloor r/(p-1)\rfloor$.

\begin{lemma}
If $p\mid N$, where $N$ is a positive integer, then
\[
 v_p\left(\int_0^1 H_{r,p}(Nx)\,dx\right)\geq\lfloor r/p\rfloor.
\]
\end{lemma}
\begin{proof}
For $\pi_h=\zeta^h-1$, define
\[
 W_{N,h}(t)=(1+\pi_h t)^N-1.
\]
Its degree-$k$ coefficient has $w$-valuation at least
$\lceil(p-1)k/p\rceil+1$.
For $k\geq p$ this follows from the factor $\pi_h^k$.
For $1\leq k<p$, the binomial coefficient $\binom Nk$ is divisible by $p$,
which supplies the additional valuation required.
Together with \eqref{eq:gj}, this shows that
\[
 F_N(\pi_h t):=\int_0^1(1+\pi_h t)^{Nx}\,dx
    =G(W_{N,h}(t))=\sum_{k\geq0}A_{k,h}t^k
\]
has $A_{0,h}=1$ and
\[
 w(A_{k,h})\geq\left\lceil\frac{p-1}{p}k\right\rceil.
\]
This is a formal identity, obtained from
$\log((1+z)^N)=N\log(1+z)$.
It can also be evaluated at $t=1$: the terms
$g_jW_{N,h}(t)^j$ have valuations at least $(p-1)j$ for integral $t$,
so their sum converges uniformly. Since $p\mid N$,
$W_{N,h}(1)=0$, and therefore $\sum_{k\geq0}A_{k,h}=1$.
Consequently,
\[
 \int_0^1H_{r,p}(Nx)\,dx
   =\frac1p\sum_{h=1}^{p-1}\sum_{k=r+1}^{\infty}A_{k,h}.
\]
The $w$-valuation is at least
$\lceil(p-1)(r+1)/p\rceil-(p-1)$.
The result is rational, so its $p$-adic valuation is at least
$\lceil(r+1)/p-1\rceil=\lfloor r/p\rfloor$.
\end{proof}

\begin{proof}[Conclusion]
Choose the smallest $a\geq1$ such that every $y_i=p^a x_i$ lies in $\Z_p$.
At least one $y_i$ is a unit. Let $q$ count these units, so $1\leq q\leq K$.
By exactness and \eqref{eq:filter},
\[
 q\equiv K\int_0^1H_{r,p}(p^a x)\,dx
             \pmod {p^{\lfloor r/(p-1)\rfloor}}.
\]
The lemma implies $p^{\lfloor r/p\rfloor}\mid q$, proving the theorem.
More precisely,
$p^{\min(\lfloor r/(p-1)\rfloor,\,v_p(K)+\lfloor r/p\rfloor)}\mid q$.
\end{proof}

\paragraph{Scope.}
No real interval restriction on the rational nodes was used. Repeated nodes
are allowed. This proves an exponential lower bound for rational interval
designs, but the one-way insertion construction in the dice manuscript does
not transfer it to arbitrary dice.
\begin{thebibliography}{9}
\bibitem{mitkin} D.~A.~Mit'kin, \emph{An estimate for the number of terms in the Hilbert--Kamke problem}, Math. USSR Sbornik \textbf{57} (1987), no.~2, 561--590, Lemma~7 and its following remark, \href{https://www.mathnet.ru/eng/sm1845}{primary journal record and full text}.
\end{thebibliography}
\end{document}
